
\emph{est dit} de type fini \emph{si toute famille limit\'ee de classes de faisceaux inversibles sur les fibres de @@M0@@ a comme image inverse une famille limit\'ee de classes de faisceaux sur les fibres de @@M1@@.}

\textbf{Lemme 3.4.} \emph{Soient @@M2@@ propres sur @@M3@@ noeth\'erien et @@M4@@ surjectif de type fini. Alors}~:
\begin{enumerate}[label=(\roman*)]
\item \emph{Pour qu'un morphisme @@M5@@ soit de type fini, il faut et il suffit que @@M6@@ soit de type fini.}
\item \emph{Dans un diagramme commutatif}
@@M7@@
\emph{avec @@M8@@ et @@M9@@ de type fini, on a @@M10@@ de type fini.}
\end{enumerate}

En effet, supposons @@M11@@ de type fini. Soit @@M12@@ une famille sur les fibres de @@M13@@, limit\'ee par un @@M14@@ sur un @@M15@@. La famille @@M16@@ sur les fibres de @@M17@@, d\'efinie de fa\c{c}on \'evidente, est limit\'ee par l'image inverse de @@M18@@ sur @@M19@@. Si @@M20@@ est limit\'e par un @@M21@@ sur un @@M22@@, alors @@M23@@ limite aussi @@M24@@.

La r\'eciproque de (i) et l'assertion (ii) se d\'emontrent de fa\c{c}on analogue.

\textbf{Th\'eor\`eme 3.5.} \emph{Soient @@M25@@ un sch\'ema noeth\'erien, @@M26@@ et @@M27@@ deux @@M28@@-sch\'emas propres et @@M29@@ un morphisme. Si @@M30@@ est surjectif, alors le morphisme induit entre les foncteurs de Picard}
@@M31@@
